\section{A uniform global envelope through the source collisions}
\label{app:global-core-envelope}

The local interpolation of Section~\ref{app:selected-source} is used
only on a fixed high-index interval. Here we prove the different estimate
needed before restricting to that interval: a uniform exponential-rate
bound for every physical color and every core index. The proof retains
the other recurrence branch, including its repeated-root colors.
Here $C_j^\eps$ always denotes the core polynomial of a single physical
leaf of color $j$ and half-twist sign $\eps$, as in
\eqref{eq:core-alias}; it does not denote the whole analytic packet.

\subsection{The complete-diagonal normalization and its literal wall}

Use $q,Q,d,M,F_r$ as in Section~\ref{app:selected-source}. For a physical
integer $0\le h\le M$, interpret the finite sum $A_{h,r}$ in
\eqref{eq:b1-A-source} with lower endpoint $\max(r,h)$. Thus all its
factorials have nonnegative arguments. Set
\[
 V_{h,r}=\frac{q^{h^2}A_{h,r}}{D_{0,r}},\qquad
 \chi=\frac{\pi h}{N},\qquad H=Q^h,\qquad Z=Q^{r+1}.
\]
The denominator is the complete diagonal, which is nonzero at every
$0\le r\le M$. The Ward identity, or its polynomial version before
cancellation of zero endpoint factors, gives
\begin{equation}\label{eq:b2-global-recurrence}
\begin{aligned}
 V_{h,r}&=p_{h,N}(r)V_{h,r+1}+q_N(r)V_{h,r+2},\\
 p_{h,N}(r)
 &=-\frac{Z^{-1}[(H+H^{-1})Z^2+1+Q^{-1}]}{(1-Z)^2},\\
 q_N(r)
 &=-\frac{Q^{-2}Z^{-2}(1-QZ^2)(1+QZ)}
 {(1-Z)^2(1-QZ)} .
\end{aligned}
\end{equation}
It holds also when $r<h$: the omitted lower terms of the polynomial
Ward telescope have their zero numerator factors, rather than divided
negative factorials. More explicitly, the certificate
$g_h(z)$ in Section~\ref{app:selected-source} vanishes at the lower
endpoint $z=Q^{\max(r,h)}$: its factor $1-z/Y$ applies when $r\ge h$,
and its factor $1-z/H$ applies when $r<h$.
At $r=M-1$ one has $q_N(r)=0$ at both parities.
We may therefore set the unused auxiliary value $V_{h,M+1}=0$.
The actual remaining terminal value obeys the exact identity
\begin{equation}\label{eq:b2-terminal}
 |V_{h,M}|=\frac{s_M^2}{s_{M-h}s_{M+h}}
 =\begin{cases}
 1,&N\ \hbox{odd},\\
 \cos\chi,&N\ \hbox{even}.
 \end{cases}
\end{equation}
Indeed $s_a s_{N-1-a}=N$. For odd $N=2M+1$, both products in
\eqref{eq:b2-terminal} equal $N$. For even $N=2m$, $M=m-1$,
$s_M^2=N/2$ and
$s_{M-h}s_{M+h}=N/(2\cos\chi)$. Thus the source at the wall has
modulus at most one uniformly over the entire physical color range.

The finite triangular transport and the marked factorial phases give
the following exact modulus in the normalized Newton basis:
\begin{equation}\label{eq:b2-source-coefficient}
 \bigl|[\Bhat_r]C_{N-h}^{\,\eps}\bigr|
   =\frac{N^2s_{2r}}{s_r^4}|V_{h,r}|,\qquad
 0\le r\le M,\quad \eps\in\{-1,1\}.
\end{equation}
One can obtain this directly by interchanging the two finite sums in
\eqref{eq:b1-signed-finite}; the coefficient in the unnormalized basis
$B_r$ is $N^2s_{2r}V_{h,r}/s_r^6$ times a unit phase.
Multiplication by $s_r^2$ gives \eqref{eq:b2-source-coefficient}.
The positive sign is the conjugate of the negative sign on the
physical grid. Phases have been used only to compute the modulus,
not deleted from the finite source.

\subsection{The spectral ellipse}

Let $T_N(r)$ be the companion matrix of
\eqref{eq:b2-global-recurrence}, so that
$(V_{h,r},V_{h,r+1})^t=T_N(r)(V_{h,r+1},V_{h,r+2})^t$.
Freeze $u=(r+1)\pi/N$ and let $Q\to1$ while
$Z=e^{2\ii u}$ and $H=e^{2\ii\chi}$ stay fixed. The limiting
matrix $T_0(u,\chi)$ has eigenvalues $Z^{-1}v$, where
\begin{equation}\label{eq:b2-ellipse-polynomial}
 v^2-bv-c=0,\qquad
 c=\cot^2u,\qquad
 b=(c-1)\cos^2\chi-2\ii\sqrt c\,\sin^2\chi.
\end{equation}
We claim
\begin{equation}\label{eq:b2-spectral-envelope}
 \rho(T_0(u,\chi))\le \rho_0(u):=\max(1,\cot^2u),
 \qquad 0<u\le\pi/2,\quad0\le\chi\le\pi/2.
\end{equation}

For $0<c<1$, a root on the unit circle would give $b=v-c/v$.
The unit-root locus in the $b$ plane is therefore the ellipse
\[
 \left(\frac{\Re b}{1-c}\right)^2+
 \left(\frac{\Im b}{1+c}\right)^2=1.
\]
For the value in \eqref{eq:b2-ellipse-polynomial}, its left side is
\[
 \cos^4\chi+\frac{4c}{(1+c)^2}\sin^4\chi\le1,
\]
strictly less than one when $\chi>0$. At $b=0$ the roots have
modulus $\sqrt c<1$. The line segment to an interior point of this
ellipse cannot cross the unit-root locus, so both roots remain in
the open unit disk. Continuity treats the ellipse boundary.
For $c>1$, apply the same argument to $v/c$, with parameter $1/c$.
At $c=1$ the roots of
$v^2+2\ii\sin^2\chi\,v-1$ have modulus one. At $c=0$ they are
$0$ and $-\cos^2\chi$. This proves
\eqref{eq:b2-spectral-envelope} without selecting an eigenbasis.

In particular, if $\chi_0>0$ and
$I\subset(\pi/4,\pi/2)$ is a fixed closed interval, then
\begin{equation}\label{eq:b2-strict-spectral}
 \sup_{\substack{u\in I\\\chi_0\le\chi\le\pi/2}}
 \rho(T_0(u,\chi))\le\theta<1.
\end{equation}
The endpoint color $\chi=\pi/2$ is included: there the two roots
coalesce, but their common modulus is $\sqrt c<1$.

\subsection{Frozen blocks and the singular lower endpoint}

Define
\begin{equation}\label{eq:b2-envelope-action}
 E(u)=
 \begin{cases}
 \displaystyle\frac2\pi\int_u^{\pi/4}\log\cot v\,\dd v,
       &0\le u<\pi/4,\\
 0,&\pi/4\le u\le\pi/2.
 \end{cases}
\end{equation}
We prove the two uniform assertions
\begin{align}
 \limsup_{N\to\infty}
 \max_{0\le h,r\le M}
 \left\{N^{-1}\log|V_{h,r}|-E(\pi r/N)\right\}
 &\le0,\label{eq:b2-global-rate}\\
 \limsup_{N\to\infty}
 \max_{\substack{\delta N\le h\le M\\0\le r\le\beta N}}
 \left\{N^{-1}\log|V_{h,r}|-E(\pi r/N)\right\}
 &\le-\gamma_{\delta,\beta},
 \label{eq:b2-far-color}
\end{align}
where $0<\delta<1/2$, $0<\beta<1/2$, and
$\gamma_{\delta,\beta}>0$. We use $\log0=-\infty$.

First fix a lower cutoff $a>0$. On $a\le u\le\pi/2$ the exact
coefficients and their first $u$ derivatives are uniformly bounded,
independently of the physical color. The exact normalized matrices
$B_N(r)=T_N(r)/\rho_0(u_r)$, $u_r=(r+1)\pi/N$, differ by
$O_a(N^{-1})$ from the bounded Lipschitz family
$B(u_r,\chi)=T_0(u_r,\chi)/\rho_0(u_r)$.

For any bounded $2$ by $2$ matrix $B$ whose eigenvalues have modulus
at most one, Cayley--Hamilton gives
\begin{equation}\label{eq:b2-CH}
 \|B^L\|\le L\|B\|+L-1\le C_aL.
\end{equation}
The coefficient of $B$ is
$\sum_{j=0}^{L-1}\lambda_1^{L-1-j}\lambda_2^j$, of modulus at
most $L$, and the constant coefficient has modulus at most $L-1$.
These expressions remain valid at a double root. If both roots
instead have modulus at most $\theta<1$, the same calculation gives
$\|B^L\|\le C_aL\theta^{L-1}$.

Choose a fixed block length $L$. Telescoping a product of $L$
adjacent exact normalized matrices against its frozen value gives
\[
 B_N(r)\cdots B_N(r+L-1)
   =B(u_r,\chi)^L+O_{a,L}(N^{-1}),
\]
uniformly in $\chi$. For a prescribed $\epsilon>0$, choose $L$
so that $C_aL<e^{\epsilon L}$, and then take $N$ sufficiently
large. Partition any propagation product into these blocks.
The incomplete end blocks cost a fixed constant, and all ordinary
blocks together cost at most $C_{a,\epsilon}e^{2\epsilon N}$.

For the strict estimate, choose
$I\subset(\max(\pi\beta,\pi/4),\pi/2)$. Every propagation from
the physical wall to $r\le\beta N$ crosses this interval. By
\eqref{eq:b2-strict-spectral}, increase the same fixed $L$ until
each full block in $I$ has norm at most $e^{-\gamma_0L}$ for some
$\gamma_0>0$. There are $N|I|/(\pi L)+O(1)$ such blocks.
Choose $\epsilon<\gamma_0|I|/(8\pi)$. Their strict negative budget
then exceeds the cost of all ordinary blocks.

Only $u=0$ has a singular coefficient. Directly from
\eqref{eq:b2-global-recurrence}, for $u_r\le a<\pi/8$,
\begin{equation}\label{eq:b2-endpoint-matrices}
 \|T_N(r)\|\le C\frac{N^2}{(r+1)^2}.
\end{equation}
Here $|H+H^{-1}|\le2$, $|1-Z|\ge4(r+1)/N$, and the ratio
$|1-QZ^2|/|1-QZ|$ is bounded by an absolute constant in this
range. If $K\le aN/\pi$, the product of the bounds in
\eqref{eq:b2-endpoint-matrices} is at most
$C^KN^{2K}/(K!)^2$. Hence the entire omitted initial segment
costs
\[
 \exp\{O(Na\log(1/a))+O(\log N)\},
\]
uniformly in its starting index. This cost has rate tending to
zero as $a\downarrow0$.

Finally $\log\rho_0(u)$ is decreasing until it reaches zero, so
the right-endpoint Riemann sum satisfies
\begin{equation}\label{eq:b2-Riemann-envelope}
 \sum_{j=r}^{M-1}\log\rho_0((j+1)\pi/N)
       \le N E(\pi r/N).
\end{equation}
Combine this scalar product with \eqref{eq:b2-terminal}, the frozen
blocks, and the endpoint bound. For \eqref{eq:b2-global-rate},
take $N\to\infty$ first and then let $a,\epsilon\downarrow0$.
For \eqref{eq:b2-far-color}, first fix $I$ and its negative
budget, then make $a\log(1/a)$ and $\epsilon$ small compared with
that budget, and finally take $N$ large. This proves both estimates
with their stated uniformity. There is no adiabatic eigenvector
argument and no polynomial bound claimed at a spectral collision.

\subsection{The actual coefficient envelopes and the macroscopic color loss}

The elementary sine-factorial estimate is uniform up to both endpoints:
\begin{equation}\label{eq:b2-factorial-endpoints}
 \log s_j=NJ(j/N)+O(\log N),\qquad 0\le j<N.
\end{equation}
For example, subtract $\log t+\log(1-t)$ from
$\log(2\sin\pi t)$ on $[0,1]$. The remainder is continuously
differentiable, so its Riemann-sum error is bounded; the two removed
terms are ordinary factorials and have $O(\log N)$ errors by
Stirling's inequalities. This also proves the uniform assertion
when one of the factorial indices is zero.

The complete prefactor in \eqref{eq:b2-source-coefficient} has rate
$g(x)=J(2x)-4J(x)$. Since $g'(x)=2\log\cot\pi x$,
\begin{equation}\label{eq:b2-g-plus-E}
 g(x)+E(\pi x)=g(\max(x,1/4))=\widetilde g(x).
\end{equation}
Consequently \eqref{eq:b2-global-rate} supplies a sequence
$e_N\downarrow0$ after taking its decreasing upper envelope, such that
\[
 |[\Bhat_r]C_{N-h}^{\,\eps}|
 \le \exp\{N[\widetilde g(r/N)+e_N]\},
 \qquad 0\le h,r\le M,\quad\eps\in\{-1,1\}.
\]
The same estimate holds for the pure sign. Its coefficient is zero
for $r<h$ and otherwise is
$N^2s_{2r}/(s_r^2s_{r-h}s_{r+h})$. Its rate is
$g(x)-\Delta(x,t)$, where
\[
 \Delta(x,t)=J(x-t)+J(x+t)-2J(x),\qquad 0\le t\le x\le1/2.
\]
For $t>0$ and $x<1/2$,
\[
 \partial_t\Delta
  =\log\frac{\sin\pi(x+t)}{\sin\pi(x-t)}>0,
\]
because the difference of the two sines is
$2\cos\pi x\sin\pi t>0$. Continuity includes the integrable
boundary $t=x$. Thus pure cores have a uniform strict loss when
$x$ lies in a fixed compact high interval and $t\ge\delta>0$.
For signed cores the same conclusion follows from
\eqref{eq:b2-far-color}, choosing $\beta$ above that high interval.
Both parities, the largest physical color, and all zero coefficients
are retained.

For completeness, the alias bound in \eqref{eq:leaf-envelopes}
requires no recurrence. We first prove the finite bound on the
untransformed alias. Write $\theta=\pi/N$. Its coefficient from
\eqref{eq:low-alias}, for $1\le j\le\lfloor N/2\rfloor$ and
$0\le r<j$, has modulus
\[
 b_{j,r}=\frac{s_r^2}{s_{2r+1}}
           \prod_{u=-r}^{r}2\sin((j+u)\theta).
\]
As a function of real $j$ in $r<j<N-r$, this positive product is
log concave and symmetric about $N/2$, so its maximum is there.
Cancelling $(2\sin u\theta)(2\cos u\theta)=2\sin2u\theta$
gives the upper bound
\[
 B_r=2\,
 \frac{\prod_{u=1}^{r}2\sin(2u\theta)}
      {\prod_{u=0}^{r}2\sin((2u+1)\theta)},\qquad
 \frac{B_{r+1}}{B_r}
    =\frac{\sin((2r+2)\theta)}{\sin((2r+3)\theta)}.
\]
The ratio is at most one and then at least one, with only one
possible switch, at $(4r+5)\theta=\pi$. Therefore the maximum
over $0\le r\le\lfloor N/2\rfloor-1$ is at an endpoint.
The first is $B_0=\csc\theta$. At the other endpoint, the even
and odd sine products give $B_{N/2-1}=N/2$ when $N$ is even;
when $N=2m+1$, reflection of the sine factors gives
$B_{m-1}=\csc\theta$. The even product identities follow by
evaluating the factorizations of $z^m-1$ and $z^m+1$ at $z=1$.
Finally $\sin(\pi/N)\ge2/N$. We have proved
$b_{j,r}\le N/2$ at every color and both parities.

A triangular half-twist entry has modulus
$s_v/(s_rs_{v-r})$, so \eqref{eq:b2-factorial-endpoints} gives
the rate
\[
 J(w)-J(u)-J(w-u),\qquad u=r/N,\quad w=v/N\ge\max(u,1/2).
\]
For $0<u<1$ its derivative in $w$ is
$\log[\sin(\pi w)/\sin(\pi(w-u))]$; it changes from positive to
negative at $w=(1+u)/2$. Reflection
$J(1-z)=-J(z)$ and duplication give the maximum
\[
 J((1+u)/2)-J(u)-J((1-u)/2)=-2J(u/2)=k(u).
\]
The two endpoint cases follow by continuity. All factorial errors
are $O(\log N)$, and there are at most $N$ summands. This proves
the alias bound with a polynomial prefactor, on its full original
index range $0\le r\le d$.

To pass from these leaf estimates to the complete source, first fix
the positive compositional alias margins and the compact core-state
gap in Chapter~\ref{ch:alias}. Every fixed tree has finitely many
leaves and gates and only polynomially many lattice states. Absorb
their uniform $e_N$ and the $O(\log N)/N$ gate errors into half the
fixed margin. Inside the remaining strict state window, the
macroscopic-color loss just proved supplies the other half of the
argument. This gives the whole-source bounds
\eqref{eq:b1-source-restoration}--\eqref{eq:b1-color-complement}.
The result is a uniform exponential-rate envelope at all colors;
the polynomial leading packet remains the separate regular local
theorem of Section~\ref{app:selected-source}.
